\section{Stress and displacement fields near the corner}
\label{app:asymptotics}
\providecommand{\SingularityExponentFigure}{Figure~\ref{fig:singularity-exponents}}

\setcounter{equation}{0}
\renewcommand{\theequation}{A.\arabic{equation}}

This appendix collects the characteristic relations and asymptotic-field
coefficients used in the main text. For plane strain, the dimensionless
elastic ratio $e$ and material parameters $\kappa_i$ are
\begin{equation*}
  e=\frac{1+\nu_2}{1+\nu_1}\frac{E_1}{E_2},
  \qquad
  \kappa_i=3-4\nu_i,\qquad i=1,2.
\end{equation*}
For each characteristic exponent, the complex Mellin variable $p$ and
the generic singularity exponent $\lambda_*$ are related by
\begin{equation*}
  p=-1-\lambda_*,
  \qquad -1<\lambda_*<0.
\end{equation*}
The physical exponent is taken to be the smallest real root in
$(-1,0)$ that is continuously connected to the corresponding solution
branch. The case $\alpha=\pi/2$ is the degenerate straight-interface
limit and is not treated as an ordinary interior point of a
characteristic branch.

\subsection{Intact corner}

For the configuration without an interfacial shear crack, the singular
part of the interfacial stress field is~\cite{NazarenkoKipnis2022Stress}
\begin{equation}
\label{eq:app-intact-field}
\begin{aligned}
  \tau_{r\theta}(r,\alpha)
    &=Cg_1r^{\lambda_0}+o\!\left(r^{\lambda_0}\right),\\
  \sigma_\theta(r,\alpha)
    &=Cg_2r^{\lambda_0}+o\!\left(r^{\lambda_0}\right),
    \qquad r\to0,
\end{aligned}
\end{equation}
where $C$ is the local singular-field amplitude, $\lambda_0$ is
the intact-corner singularity exponent, and the dimensionless angular
coefficients $g_1$ and $g_2$ are
\begin{equation*}
\begin{aligned}
  g_1={}&(1-e)\sin[\lambda_0(\pi-\alpha)]
  \bigl[(\lambda_0+1)\sin2\alpha
       +\sin\{2(\lambda_0+1)\alpha\}\bigr]\\
  &+e(1+\kappa_2)\sin(\lambda_0\pi)
      \sin[(\lambda_0+2)\alpha],\\[2pt]
  g_2={}&(1-e)\cos[\lambda_0(\pi-\alpha)]
  \bigl[(\lambda_0+1)\sin2\alpha
       +\sin\{2(\lambda_0+1)\alpha\}\bigr]\\
  &+e(1+\kappa_2)\sin(\lambda_0\pi)
      \cos[(\lambda_0+2)\alpha].
\end{aligned}
\end{equation*}
The exponent $\lambda_0$ is determined by the characteristic equation
\begin{equation}
\label{eq:app-D0-root}
  D_0(-1-\lambda_0)=0,
  \qquad -1<\lambda_0<0,
\end{equation}
where the intact-corner determinant $D_0(p)$ is
\begin{equation*}
  D_0(p)=b_0(p)+e\,b_1(p)+e^2b_2(p),
\end{equation*}
with auxiliary functions $b_0(p)$, $b_1(p)$, and $b_2(p)$ given by
\begin{equation*}
\begin{aligned}
  b_0(p)={}&
  [\sin(2p\alpha)+p\sin(2\alpha)]
  [\kappa_1\sin\{2p(\pi-\alpha)\}+p\sin(2\alpha)],\\[2pt]
  b_1(p)={}&(1+\kappa_1)(1+\kappa_2)\sin^2(p\pi)\\
  &-[\sin(2p\alpha)+p\sin(2\alpha)]
    [\kappa_1\sin\{2p(\pi-\alpha)\}+p\sin(2\alpha)]\\
  &-[\sin\{2p(\pi-\alpha)\}-p\sin(2\alpha)]
    [\kappa_2\sin(2p\alpha)-p\sin(2\alpha)],\\[2pt]
  b_2(p)={}&
  [\sin\{2p(\pi-\alpha)\}-p\sin(2\alpha)]
  [\kappa_2\sin(2p\alpha)-p\sin(2\alpha)].
\end{aligned}
\end{equation*}
\subsection{Corner with a stationary V-shaped interfacial shear crack}

For the prescribed stationary V-shaped interfacial shear crack, the two
symmetric branches have length $l$ each. At distances $r\ll l$, the normal stress along the bisector
of the sector occupied by material~$i$ has the asymptotic
form~\cite{NazarenkoKipnis2022Stress}
\begin{equation}
\label{eq:app-cracked-field}
  \sigma_\theta(r,\beta_i)
  =CQ_i r^\lambda+o\!\left(r^\lambda\right),
  \qquad r\to0,
  \qquad
  \beta_1=\pi,\quad \beta_2=0.
\end{equation}
Here, $\lambda$ is the corner singularity exponent for the cracked
configuration and $\beta_i$ is the corresponding bisector direction.
The amplitude coefficient $Q_i$ is expressed as
\begin{equation*}
  Q_i=g_1q_i l^{\lambda_0-\lambda},
\end{equation*}
where the dimensionless coefficient $q_i$ and the derivative $D'(p)$
of the cracked-corner determinant are
\begin{equation*}
  q_i=
  \frac{2S_i(-1-\lambda)K^+(-1-\lambda)G^+(-1-\lambda)}
  {(\lambda-\lambda_0)D'(-1-\lambda)
   K^+(-1-\lambda_0)G^+(-1-\lambda_0)},
  \qquad
  D'(p)=\frac{\dd D(p)}{\dd p}.
\end{equation*}
The functions $S_i(p)$ are auxiliary numerator functions defined below.
The factors $K^+(p)$ and $G^+(p)$ are specified later in this subsection.
The base plus factor $G^+(p)$ must not be confused with $G_i^+(p)$, which
factorizes the process-zone kernel $G_i(p)$ in
Section~\ref{sec:derivation}. Explicitly,
\begin{equation*}
\begin{aligned}
  S_1(p)={}&2e(1+\kappa_2)s_{11}(p)\sin[(p+1)\alpha]\\
  &+s_{12}(p)
  \Bigl[(1+\kappa_1)\sin\{p(\pi-\alpha)-\alpha\}
        +2(1-e)s_{13}(p)\Bigr],\\[2pt]
  S_2(p)={}&s_{21}(p)
  \Bigl[e(1+\kappa_2)\sin[(p+1)\alpha]
        +2(1-e)s_{22}(p)\Bigr]\\
  &-2(1+\kappa_1)\sin\{p(\pi-\alpha)-\alpha\}\,s_{23}(p),
\end{aligned}
\end{equation*}
where the auxiliary functions $s_{mn}(p)$, with $m=1,2$ and
$n=1,2,3$, are
\begin{equation*}
\begin{aligned}
  s_{11}(p)&=p\cos(p\pi-\alpha)\sin\alpha
             -\sin[p(\pi-\alpha)]\cos(p\alpha),\\
  s_{12}(p)&=p\sin2\alpha+\sin(2p\alpha),\\
  s_{13}(p)&=p\cos[p(\pi-\alpha)]\sin\alpha
             -\sin[p(\pi-\alpha)]\cos\alpha,\\
  s_{21}(p)&=\sin[2p(\pi-\alpha)]-p\sin2\alpha,\\
  s_{22}(p)&=p\cos(p\alpha)\sin\alpha
             +\sin(p\alpha)\cos\alpha,\\
  s_{23}(p)&=p\cos(p\pi+\alpha)\sin\alpha
             +\cos[p(\pi-\alpha)]\sin(p\alpha).
\end{aligned}
\end{equation*}
The corner singularity exponent $\lambda$ for the configuration
containing the prescribed stationary V-shaped interfacial crack satisfies
\begin{equation}
\label{eq:app-D-root}
  D(-1-\lambda)=0,
  \qquad -1<\lambda<0,
\end{equation}
where the characteristic determinant $D(p)$ is
\begin{equation*}
  D(p)=a_0(p)+e\,a_1(p),
\end{equation*}
with auxiliary functions $a_0(p)$ and $a_1(p)$ given by
\begin{equation*}
\begin{aligned}
  a_0(p)={}&(1+\kappa_1)
  [\cos\{2p(\pi-\alpha)\}-\cos2\alpha]
  [\sin(2p\alpha)+p\sin2\alpha],\\
  a_1(p)={}&(1+\kappa_2)
  [\cos(2p\alpha)-\cos2\alpha]
  [\sin\{2p(\pi-\alpha)\}-p\sin2\alpha].
\end{aligned}
\end{equation*}
The coefficients $q_i$ use the base kernel
\begin{equation*}
  G(p)=
  \frac{(e+\kappa_1)(1+e\kappa_2)\sin(\pi p)D(p)}
  {(e+\kappa_1+1+e\kappa_2)\cos(\pi p)D_0(p)}
\end{equation*}
and its plus factor
\begin{equation*}
  G^+(p)=
  \exp\!\left\{
    \frac{1}{2\pi\ii}
    \int_{-\ii\infty}^{\ii\infty}
    \frac{\log G(z)}{z-p}\,\dd z
  \right\},
  \qquad \operatorname{Re}p<0.
\end{equation*}
Here, $z$ is the complex integration variable, $\ii^2=-1$, and the contour
is directed upward along the imaginary axis. The gamma-function factor is
\begin{equation*}
  K^+(p)=\frac{\Gamma(1-p)}{\Gamma(\tfrac12-p)}.
\end{equation*}
Here, $\Gamma$ is the Euler gamma function. For real $p<0$, $G^+(p)$
can equivalently be evaluated using the real integral
\begin{equation*}
  \log G^+(p)
  =-\frac{p}{\pi}\int_0^\infty
    \frac{\log G(\ii t)}{t^2+p^2}\,\dd t.
\end{equation*}
Here, $t$ is a real integration variable.

\subsection{Corner with a process zone}

For a process zone in material~$i$, the kernel of the functional equation
in Section~\ref{sec:derivation} is
\begin{equation*}
  G_i(p)=\frac{D_i(p)\cos(\pi p)}{D(p)\sin(\pi p)}.
\end{equation*}
The corresponding characteristic determinants are
\begin{equation}
\label{eq:app-post-zone-determinants}
\begin{aligned}
  D_1(p)={}&2e(1+\kappa_2)
  [\cos(2p\alpha)-\cos2\alpha]
  \{\sin^2[p(\pi-\alpha)]-p^2\sin^2\alpha\}\\
  &+(1+\kappa_1)
  [\sin(2p\alpha)+p\sin2\alpha]
  [\sin\{2p(\pi-\alpha)\}-p\sin2\alpha],\\[3pt]
  D_2(p)={}&2(1+\kappa_1)
  [\cos\{2p(\pi-\alpha)\}-\cos2\alpha]
  [\sin^2(p\alpha)-p^2\sin^2\alpha]\\
  &+e(1+\kappa_2)
  [\sin\{2p(\pi-\alpha)\}-p\sin2\alpha]
  [\sin(2p\alpha)+p\sin2\alpha].
\end{aligned}
\end{equation}
% The author-confirmed signs in D_1 are preserved. Their provenance is
% recorded in EDITORIAL_LOG.md rather than in the clean manuscript.
The exponent $\lambda_i$ for a configuration with a zone in
material~$i$ is determined by
\begin{equation}
\label{eq:app-post-zone-root}
  D_i(-1-\lambda_i)=0,
  \qquad -1<\lambda_i<0,
  \qquad i=1,2.
\end{equation}
These roots define the $\lambda_1$ and $\lambda_2$ curves in
\SingularityExponentFigure. They characterize the corner singularity in the
presence of a process zone but are not used in the formulas for $d_i$,
$\delta_i$, and $J_i$; those formulas use $\lambda$ from
Eq.~\eqref{eq:app-D-root}.
